\documentclass[letterpaper]{article}
\PassOptionsToPackage{unicode=true}{hyperref} % options for packages loaded elsewhere
\PassOptionsToPackage{hyphens}{url}
\usepackage{graphicx} % Required for inserting images


\def\tightlist{}


\graphicspath{ {../images/} }
\usepackage[margin=1in]{geometry}
\usepackage{tikz}
\usepackage[dvipsnames]{xcolor}
\usetikzlibrary{calc}




\usepackage[utf8]{inputenc}
\usepackage[T1]{fontenc}
\usepackage{lmodern}
\usepackage[english]{babel}

\usepackage{longtable}
\usepackage{booktabs}

% Useful packages
\usepackage{hyperref}
\usepackage{amsmath, amssymb}
\usepackage{fancyhdr}
\usepackage{setspace}
\usepackage{tocloft}
\usepackage{titlesec}


\usepackage{unicode-math}



% Code highlighting (Pandoc uses listings or minted, but listings is safer without shell escape)
\usepackage{listings}
\lstset{
  basicstyle=\ttfamily\small,
  breaklines=true,
  frame=single,
  backgroundcolor=\color[gray]{0.95}
}

\usepackage{setspace}


%Modifiable Commands; defaults set to none
\newcommand{\titleDOC}{}
\newcommand{\authorDOC}{Daniel Sabalakov}
\newcommand{\dateDOC}{\today}
\newcommand{\classDOC}{}
\newcommand{\emailDOC}{danielsabalakov@gmail.com}
% DEFAULT QUOTE IS BY ISSAC NEWTON. IF YOU DON'T WANT IT, DO: quote: @none
\newcommand{\quoteDOC}{``\textit{Truth is ever to be found in the simplicity, and not in the multiplicity and confusion of things}'' - Issac Newton}

% Header/Footer
\pagestyle{fancy}
\fancyhf{}
\rhead{\thepage}
\lhead{\leftmark}
\fancyhead[R]{\thepage} % Place page number in top right corner




% %Title Arguments
% \title{\TITLEOFDOCUMENT}
% \author{\AUTHOROFDOCUMENT}
% \date{\today}


%
% ANYTHING BELOW HERE CAN GET PROCEDURALLY GENERATED
%
% BEGINNING OF DOCUMENT
%
%
%
\begin{document}

%titlepage








\renewcommand{\titleDOC}{2.1 Magazine Article}

\renewcommand{\authorDOC}{Daniel Sabalakov}

\renewcommand{\quoteDOC}{``\textit{May the Force Be With You}'' \\ \textemdash{} General Jan Dodonna}

\renewcommand{\classDOC}{PE Running}





\begin{titlepage}
  \titleDOC
  \classDOC
  \dateDOC
  \authorDOC
\end{titlepage}

\newpage

\section{Back to the Basics... Principles of Exercise Revisited}\label{back-to-the-basics...-principles-of-exercise-revisited}

Warming up and cooling down is an integral part of any exercise, but especially running. Failure to warm up or cool down may lead to soreness, fatigue, or injury. Warming up increases the
temperature of the body, making muscles more flexible. More flexible muscles means that the muscle is less likely to be torn or pulled, both detrimental injuries. Warming up allows the body to slowly adjust to an elevated form of stress, which reduces the buildup of lactic acid. According to active.com, "a good and consistent stretching program can save you a lot of trouble and keep you running when you might otherwise become injured." Cooling down helps adjust the body slowly adjust to a period of lower stress. When your heart rate is elevated after intense physical activity, stopping rapidly makes the blood flow decrease rapidly. This makes your blood pool at your feet. Without a period to cool down, muscles may tighten up and become sore.

Now that we know the importance of warming up and cooling down, lets talk about some tips. When you warm up and cool down, you must not rush through it. Warming up and cooling down is a delicate process; this following list of guidelines summarizes the "do\textquotesingle s" and "don\textquotesingle ts" of stretching before a run.

\begin{enumerate}
\def\labelenumi{\arabic{enumi}.}
\tightlist
\item
  Do not bounce.

  \begin{itemize}
  \tightlist
  \item
    Bouncing may cause injury because it can risk pulling or tearing a muscle. If the stretch is done too fast, it may shock the muscle into contracting quickly.
  \end{itemize}
\item
  Build a regular schedule around training

  \begin{itemize}
  \tightlist
  \item
    Most of the time, it is better for your body to have a shorter workout so that you can stretch. Stretch before and after your run to allow your body to have the benefits of warming up and cooling down.
  \end{itemize}
\item
  Avoid static stretches immediately

  \begin{itemize}
  \tightlist
  \item
    Begin with dynamic stretches to allow your muscle to warm up, and then you can lean into static stretches. Beginning static stretches immediately may risk injury.
  \end{itemize}
\end{enumerate}

The following exercises have been compiled for runners:

\begin{enumerate}
\def\labelenumi{\arabic{enumi}.}
\tightlist
\item
  Jog 3-5 minutes with various dynamic dynamic stretches in between. This fulfills the dynamic portion of warming up.

  \begin{itemize}
  \tightlist
  \item
    Exercises such as high knees and butt kicks allow all parts of the body to be warmed up.
  \item
    Other dynamic exercises include leg swings, walking lunges, and skips.
  \end{itemize}
\item
  Incorporate static stretches

  \begin{itemize}
  \tightlist
  \item
    Hold various stretches for 20-30 seconds. This should not be painful. Stretches include hamstring stretches, calf stretches, and hip flexor stretches.
  \end{itemize}
\end{enumerate}

To run efficiently, runners generally should keep a proper posture. Shoulders should be back and relaxed with no arch in the lower back. A runner\textquotesingle s abdomen should be tight. Avoid technique, and try to strike your foot below your shoulders. Run at about a steady 180 beats per minute. Running technique and posture is important to, again, prevent injury. Potential over-striding or under-striding may cause injuries or stresses on your muscles.

\section{Citation (links)}\label{citation-links}

\begin{itemize}
\tightlist
\item
  https://www.active.com/running/articles/stay-loose-stretches-for-runners
\end{itemize}

\newpage



\end{document}
